\honest{The Third Alternative}{Eliezer Yudkowsky}{2007}

I open with an argument, in quotation marks, that I attribute to no one: believing in Santa Claus gives children a sense of wonder and makes them behave, so it is ``a Noble Lie whose net benefit should be preserved.''

This, I say, is a false dilemma. Even if believing in Santa beats doing nothing, it need not be the best choice. A Space Shuttle launch or science fiction can supply wonder. And bribes make children behave ``only'' when adults are watching, ``while praise without bribes leads to unconditional good behavior.'' \nb{No source is given. A 1999 meta-analysis of 128 experiments found that tangible rewards reduced later interest in a task and that praise increased it, though praise tended to help children less than college students. It measured interest in tasks, not good behavior, and found average effects, not ``unconditional'' ones.}

From this one example I conclude that Noble Lies are ``generally'' package deals, and that if we need the benefit, ``we can construct a Third Alternative for getting it.''

The hard steps, I say, are deciding to look for an alternative and deciding to accept it, and ``most people fail on these two steps.'' I give no evidence. Some false dilemmas are honest. Others come from defending a policy by its benefit over doing nothing, and then the defender ``does not want'' a better option. ``The best is the enemy of the good,'' I say, meaning that a better option is the enemy of such a defense. \nb{The proverb usually means the reverse: holding out for the best can block a good improvement.}

Modern cognitive psychology, I say, views decisions as searches, and a search needs a rule for when to stop: you cannot look at every house in the city. But we stop at once when we find something defensible and convenient, and keep looking when the options are costly to us. \nb{The post gives no evidence, but experiments by Ditto and Lopez (1992) found something close: people needed less information to reach a conclusion they preferred. Yudkowsky returned to the idea in ``Motivated Stopping and Motivated Continuation.''} So beware of calling a policy ``defensible'' rather than ``optimal.'' Noble Liars, I add, keep whatever lie they started with. I name none.

My first test: would you be glad to find a better option, or feel ``a tiny flash of reluctance''? \nb{Three paragraphs earlier the post called the motives at work ``subconscious influences.'' It does not say why asking yourself would reveal them.} My second: have you spent five minutes, ``by the clock,'' with your eyes closed, looking for alternatives? And were you careful not to think of a good one?

I close: ``It's amazing how many Noble Liars and their ilk are eager to embrace ethical violations'' without five minutes of looking.
